% ----------------------------------------------------------
% Seção 2 — Estrutura do trabalho acadêmico (NBR 14724)
% ----------------------------------------------------------
\chapter{Estrutura do Trabalho Acadêmico}\label{cap:estrutura}

A NBR 14724 \cite{abnt14724} define a estrutura dos trabalhos acadêmicos (teses, dissertações e trabalhos de conclusão de curso) em três partes: elementos pré-textuais, textuais e pós-textuais. O \autoref{qua:estrutura} lista os elementos na ordem em que aparecem e indica quais são obrigatórios.

\begin{quadro}[htbp]
	\caption{Elementos do trabalho acadêmico}\label{qua:estrutura}
	\centering
	\footnotesize
	\begin{tabular}{p{0.12\textwidth}p{0.33\textwidth}p{0.13\textwidth}p{0.30\textwidth}}
		\toprule
		\textbf{Parte} & \textbf{Elemento} & \textbf{Uso} & \textbf{Comando} \\
		\midrule
		Externa & Capa & obrigatório & \verb|\imprimircapa| \\
		& Lombada & opcional & --- \\
		\midrule
		Pré-textual & Folha de rosto & obrigatório & \verb|\imprimirfolhaderosto| \\
		& Errata & opcional & \texttt{errata} \\
		& Folha de aprovação & obrigatório & \verb|\imprimirfolhadeaprovacao| \\
		& Dedicatória & opcional & \texttt{dedicatoria} \\
		& Agradecimentos & opcional & \texttt{agradecimentos} \\
		& Epígrafe & opcional & \texttt{epigrafe} \\
		& Resumo na língua vernácula & obrigatório & \texttt{resumo} \\
		& Resumo em língua estrangeira & obrigatório & \texttt{resumo} \\
		& Lista de ilustrações & opcional & \verb|\listoffigures*| \\
		& Lista de tabelas & opcional & \verb|\listoftables*| \\
		& Lista de abreviaturas e siglas & opcional & \texttt{siglas} \\
		& Lista de símbolos & opcional & \texttt{simbolos} \\
		& Sumário & obrigatório & \verb|\tableofcontents*| \\
		\midrule
		Textual & Introdução, desenvolvimento e conclusão & obrigatório & \verb|\chapter| \\
		\midrule
		Pós-textual & Referências & obrigatório & \verb|\imprimirreferencias| \\
		& Glossário & opcional & --- \\
		& Apêndice & opcional & \texttt{apendicesenv} \\
		& Anexo & opcional & \texttt{anexosenv} \\
		& Índice & opcional & --- \\
		\bottomrule
	\end{tabular}
	\fonte{Adaptado de \citeonline{abnt14724}.}
\end{quadro}

Os elementos pré-textuais e pós-textuais estão no \texttt{main.tex}. Os opcionais vêm desativados entre \verb|\iffalse| e \verb|\fi|; para usar um deles, apague essas duas linhas.

\section{Elementos pré-textuais}\label{sec:pre-textuais}

A capa, a folha de rosto e a folha de aprovação são geradas a partir dos dados informados no preâmbulo: \verb|\titulo|, \verb|\autor|, \verb|\orientador|, \verb|\instituicao|, \verb|\campus|, \verb|\curso|, \verb|\local|, \verb|\data| e \verb|\preambulo|. Os membros da banca são informados com \verb|\membrobancaA| e \verb|\membrobancaB|, e a sigla da instituição de cada um vai entre colchetes, por exemplo \verb|\membrobancaB[UFMG]{Nome}|.

O resumo segue a NBR 6028 \cite{abnt6028}: parágrafo único, com 150 a 500 palavras em trabalhos acadêmicos, seguido das palavras-chave, separadas entre si por ponto e vírgula e finalizadas por ponto. O resumo em língua estrangeira é a tradução do resumo em português.

As listas de ilustrações, de tabelas, de siglas e de símbolos e o sumário \cite{abnt6027} são montados automaticamente a partir do texto. Cada sigla e cada símbolo é declarado uma única vez no \texttt{main.tex}.

\section{Elementos textuais}\label{sec:textuais}

A parte textual tem introdução, desenvolvimento e conclusão. Desde a edição de 2024, a NBR 14724 determina que o trabalho seja organizado em seções, e não em capítulos. A divisão do desenvolvimento em seções primárias depende do tipo de trabalho e da área; uma divisão comum em trabalhos de conclusão de curso é: referencial teórico, trabalhos relacionados, metodologia, resultados e discussão.

\section{Elementos pós-textuais}\label{sec:pos-textuais}

As referências são geradas a partir do arquivo \texttt{referencias.bib}, como descrito na \autoref{cap:referencias}. Apêndices são textos elaborados pelo próprio autor; anexos são documentos de outros autores. Os dois são identificados por letras maiúsculas consecutivas, como mostram o \autoref{apendice:questionario} e o Anexo~\ref{anexo:lei}.

\section{Regras gerais de formatação}\label{sec:formatacao}

A classe abn\TeX2 e o pacote da UFV já aplicam as regras de apresentação da NBR 14724. Elas estão resumidas na \autoref{tab:formatacao} para que você saiba o que não deve alterar.

\begin{table}[htbp]
	\IBGEtab{%
		\caption{Regras de formatação aplicadas pelo modelo}\label{tab:formatacao}
	}{%
		\begin{tabular}{ll}
			\toprule
			Item & Regra \\
			\midrule
			Papel & A4 (\SI{21}{\centi\metre} $\times$ \SI{29,7}{\centi\metre}) \\
			Margens (anverso) & esquerda e superior: \SI{3}{\centi\metre}; direita e inferior: \SI{2}{\centi\metre} \\
			Fonte do texto & tamanho 12 \\
			Fonte menor & citações longas, notas, legendas e fontes \\
			Espaçamento & 1,5 entre linhas; simples em citações longas, notas e referências \\
			Paginação & contada a partir da folha de rosto e impressa a partir da introdução \\
			Seções & numeração progressiva até a seção quinária \\
			\bottomrule
		\end{tabular}%
	}{%
		\fonte{\citeonline{abnt14724} e \citeonline{abnt6024}.}%
	}
\end{table}

\section{Numeração progressiva}\label{sec:numeracao}

A NBR 6024 \cite{abnt6024} organiza as seções em até cinco níveis, numerados em algarismos arábicos, com destaque tipográfico gradativo. Os comandos correspondentes são \verb|\chapter| (seção primária), \verb|\section| (secundária), \verb|\subsection| (terciária), \verb|\subsubsection| (quaternária) e \verb|\subsubsubsection| (quinária). Cada seção primária começa em nova página (em página ímpar, na impressão em frente e verso). As subdivisões a seguir mostram os níveis restantes.

\subsection{Seção terciária}\label{sec:terciaria}

Texto da seção terciária.

\subsubsection{Seção quaternária}\label{sec:quaternaria}

Texto da seção quaternária.

\subsubsubsection{Seção quinária}\label{sec:quinaria}

Texto da seção quinária. Não use níveis abaixo deste; se precisar enumerar mais subdivisões, use alíneas, apresentadas na \autoref{sec:alineas}.
